\chapter{Les carnets de Daniel Quillen}\label{ch:quillen}

\section{La source}

Le Clay Mathematics Institute a mis en ligne les carnets de recherche de Daniel Quillen,
de 1968 à 2003, avec un index établi par Glenys Luke
(\url{https://www.claymath.org/online-resources/quillen-notebooks/}). Le relevé du projet
compte 630 PDF, environ 6,4~Go et 29\,000 pages, en comptant des notes de cours et un
dossier de pièces diverses. Les carnets sont en anglais, tenus comme un journal daté.

\section{Ce qui change}

\paragraph{L'accès.} Le serveur de la Clay laisse afficher ses PDF dans une autre page : le
site les montre directement, sans relais. Le nombre de pages de chaque PDF a été lu dans la
structure du fichier par des requêtes partielles, sans télécharger les volumes.

\paragraph{Les pages.} Une page de PDF montre parfois deux pages de carnet ; certaines pièces
sont des dactylographies d'autres personnes. Les dates de Quillen, en tête de ses entrées,
font de ces carnets le fonds le mieux daté de ce livre.

\section{Les droits, et ce qu'on a décidé}

Quillen est mort en 2011. Les droits sont ceux de sa famille, les numérisations celles de
la Clay, dont la page ne porte qu'un « © Clay ». Le projet en a tiré les conclusions les
plus strictes :
\begin{itemize}
  \item les transcriptions ne sont jamais versionnées dans le dépôt public ; elles sont
        dans un dépôt privé ;
  \item le site ne les affiche pas ;
  \item depuis la fin de septembre 2026, on n'en transcrit plus ;
  \item l'article Lean énonce tout dans ses propres mots et ne cite que le carnet et la
        page.
\end{itemize}

\section{Ce qui a été produit}

L'article sur les carnets démontre en Lean dix énoncés tirés des carnets de 1968 à 1983, en topologie
générale, en théorie des groupes, en algèbre linéaire et en combinatoire. Trois tiennent
sous des hypothèses plus faibles que celles du carnet, et quatre sont des pistes qu'aucune
recherche n'a retrouvées dans la littérature.

\begin{exemple}
Un fonds sous droits, bien servi, se travaille quand même : on n'y publie pas la
transcription, on publie ce qu'on en a tiré, dans ses mots, avec la référence exacte. Le
lecteur ouvre la page chez la Clay et vérifie.
\end{exemple}
